% Source: https://artofproblemsolving.com/community/c3h3126960_combination_problem
% Stable identity: aops-topic-3126960-q01
(a) Modulo 5, $\binom n2=\dfrac{n(n-1)}{2}$ can be only 0,1,3, as checking $n\equiv0,1,2,3,4$ gives 0,0,1,3,1. But $1434\equiv4\pmod5$, so $\boxed{\text{none}}$ exist. (b) Put $M=10000C!$. The integer polynomial $F(n)=n(n-1)\cdots(n-C+1)$ satisfies $F(n+M)\equiv F(n)\pmod M$. Dividing its difference by $C!$ shows $\binom{n+M}{C}\equiv\binom nC\pmod{10000}$. Test every $n=C,C+1,\ldots,C+M-1$ and retain those with residue 1434. Every solution is one retained value plus a nonnegative multiple of $M$, by division into these consecutive blocks. Thus there are infinitely many exactly when the finite list is nonempty; otherwise there are none. This gives a complete, though not optimized, algorithm for every $C$.
